\documentclass[10pt]{article}

\usepackage[utf8]{inputenc}

\usepackage[T1]{fontenc}

\usepackage{amsmath}

\usepackage{amsfonts}

\usepackage{amssymb}

\usepackage[version=4]{mhchem}

\usepackage{stmaryrd}

\usepackage{bbold}

\usepackage{graphicx}

\usepackage[export]{adjustbox}

\graphicspath{ {./images/} }



\begin{document}

спектра д. о., порожденного выражением $l u=-u^{\prime \prime}+$ $+q(x) u$ в $L_{2}(-\infty . \infty)$, необходимо и достаточно, чтобы для любого $j>0$



\[

\lim _{|x| \rightarrow \infty} \int_{x}^{x+j} q(t) d t=\infty .

\]



Для д. о. с частными производными обобщение этого критерия принимает более сложный вид. Имеются другие, более простые признаки дискретности спектра д. о., вапр. самосопряженный д. о., порожденный выражением (5), будет иметь дискретный спектр, если $q(x) \rightarrow \infty$ при $|x| \rightarrow \infty$; самосопряженный д. о. $L_{m}$ имеет дискретный спектр.



Изучение природы спектра при наличии непрерывной части становится трудной задачей. Вот нек-рые результаты: 1) если обыкновенный д. о. определяется формально самосопряженным выражением (1) с периодич. коэффициентами на $(-\infty, \infty)$ с общим периодом, то его спектр непрерывный и состоит из последовательности непересекающихся интервалов, ковцы к-рых стремятся к $-\infty$ или $+\infty ; 2)$ если д. о. определяется выражением $(-1)^{k}\left(D_{1}^{2}+\ldots+D_{n}^{2}\right)^{k}+q(x) B L_{2}\left(\mathbb{R}^{n}\right) \quad$ и $\quad \lim q(x)=$ $=0, \quad|x| \rightarrow+\infty$, непрерывный спектр его заполняет $[0, \infty]$, а отрицательный спектр дискретный и может иметь предельную точку в нуле. Если $k=1,|q(x)| \leqslant$ $\leqslant M(|x|)$ и



\[

\int^{\infty} r M(r) d r<\infty \quad\left(M(r)=O\left(r^{-1}\right)\right),

\]



то отрицательный спектр будет конечным (на непрерывном спектре нет собственных значений).



Природа спектра зависит также от граничных условий. В ограниченной области описаны конкретные граничные условия, при выполнении к-рых самосопряженный д. о. Лапласа имеет вепрерывную часть спектра. Это - результат бесконечности индексов дефекта минимального д. о. Лапласа в области с границей.



\includegraphics[max width=\textwidth, center]{2023_07_09_cdda5a89673c2d885694g-1}

изучаются с делью решения смешанных задач для дифферевциальных уравнений, а также для внутренних задач теории д. о. Пусть $L$ - эллиптич. Д. о. порядка $m$. Хорошо изучена резольвента $(L+\lambda)^{-1}$ при $\lambda>0$ и функции $\exp (-L t)$ и $\exp \left(i L^{1 / m} t\right)$ при $t>0$. Последние являются разрешающими операторами для обобщенного уравнения теплопроводности $u_{t}=-L u, u(0, x)=f(x)$ и обобщенного волнового уравнения $u_{t}=i L^{1 / m} u, u(0, x)=$ $=f(x)$. Все три оператор-функции являются интегральными с ядрами $R(x, y, \lambda), K(x, y, t), G(x, y, t)$ (функции Грина) соответственно. Формула



\[

R(x, y, \lambda)=\int_{0}^{\infty} e^{-\lambda t} K(x, y, t) d t

\]



устанавливает связь между $R$ и $K$. Нек-рые свойства $R(x, y, \lambda)$ : если $L-$ эллиптический самосопряженный д. о. порядка $m$ в $L_{2}\left(\Omega_{n}\right)$, то при $p>n / 2 m$ оператору $(L+\lambda)^{-p}$ отвечает ядро типа Карлемана; при $p>n / m$ оператор $\left(L_{\boldsymbol{m}}+\lambda\right)-p$ является ядерным и поэтому



\[

S_{p}(L+\lambda)^{-p}=\sum_{k=1}^{\infty}\left(\lambda_{k}+\lambda\right)^{-p}

\]



где $\left\{\lambda_{k}\right\}-$ собственные значения д. о. $L_{m}$. Имеются и другие признаки ядерности оператора $(L+\lambda)^{-p}$ в $L_{2}\left(\mathbb{R}^{n}\right)$.



Аналитич. и асимптотич. свойства функций Грина дают полезную информацию о спектральных характеристиках д. о. $L$. Напр., если в (13) известно поведение $S_{p}(L+\lambda)-p$ при $\lambda \rightarrow \infty$, то применение тауберовых теорем позволяет найти асимптотику $\lambda_{k}$. То же самое удается, если известна асимптотика $S_{p} \exp (-L t)$ при $t \rightarrow+0$. Асимптотика $R(x, y, \lambda)$ и $K(x, y, t)$ устанавливается, напр., методом параметрикс, методом потенциалов и т. д. Так, на этом пути удалось наїіти асимптотику $\lambda_{k}$ для обширного класса эллиштич. д. о. Для определения асимштотики спектрального ядра $E(x, y, \lambda)$ әллиптич. д. о. оказалось әффективным изучение аспмптотики ядра $G(x, y, t)$ при $t \rightarrow 0$ с дальнейппим применением различных тауберовых теорем. В частности, при $x=y$, $x \notin \Gamma$



\[

E(x, x, \lambda)=(2 \pi)^{-n} \int_{l_{0}(x, \xi)<\lambda} d \xi+O\left(\lambda^{(n-1) / m}\right) .

\]



Н е с а м о с о п р я ж е в н ы е д. о. Наиболее полные результаты имеются для обыкновенных д. о. ва конечном отрезке. Пусть д. о. $L$ определяется выражением (1) при $n=1, a_{m}(x) \equiv 1$ на функциях, имеющих $(m-1)$ абсолютно непрерывные пгроизводные и удовлетворяющие краевым условиям:



\begin{center}

\includegraphics[max width=\textwidth]{2023_07_09_cdda5a89673c2d885694g-1(1)}

\end{center}



\[

\begin{aligned}

& +\beta_{v} u^{\left(k_{v}\right)}(1)+\sum_{j=0}^{k_{v}-1} \beta_{v_{j}} u^{(j)}(1)=0, \quad 1 \leqslant v \leqslant m .

\end{aligned}

\]



Здесь $m-1 \geqslant k_{1} \geqslant \ldots \geqslant k_{m} \geqslant 0, k_{v-2}<k_{v}$ и числа $a_{v}$, $b_{v}$ одновременно не равны нулю. Пусть краевые условия (2) являются регулярными. Таковыми будут краевые условия типа Штурма - Лиувилля $\left(m=2 k, l_{v_{0}}(u)=\right.$ $\left.=l_{v_{1}}(u)=0,1 \leqslant v \leqslant m-1\right)$ и периодич. типа $\left(\alpha_{v}=\beta_{v}=1\right)$. Тогда д. о. $L$ имеет бесконечное число собственных значений, к-рые имеют точную асимптотику; система собственных и присоединенных функций д. о. $L$ полна в $L_{p}(0,1)$; разложение функций $f(x)$ из $D(L)$ по собственным и присоединенным функциям д. о. $L$ сходится равномерно на $(0,1]$. Факт полноты системы собственных и присоединенных функций имеет место также при нек-рых нерегулярных краевых условиях, в частности типа распадающихся $\left(l_{v_{0}}(u)=0, \quad 1 \leqslant v \leqslant m_{1}, \quad l_{v_{1}}(u)=0\right.$, $\left.1 \leqslant v \leqslant m_{2}, \quad m_{1} \neq m_{2}, \quad m_{1}+m_{2}=m\right)$, однако сходимость разложения в ряд по собственным и присоединенным функциям справедлива только для узкого класса (lаналитических) функций.



Пусть $L_{0}$ - самосопряженный оператор в сепарабельном гильбертовом пространстве $H$ - имеет собственные значения $\left\{\lambda_{k}\right\}$ и при нек-ром $p>0$ оператор $L_{0}^{-p}$ является ядерным. Пусть $L_{1}-$ другой оператор такой, что $L_{1} L_{0}^{-1}$ является вполне непрерывным. Тогда система собственных и присоединенных векторов оператора $L_{0}-L_{1}$ полна в $H$ (т е o p е м а $К$ е л ды ша). Применение этой теоремы дает классы д. о., к-рые имеют полную систему собственных и присоединенных функций.



Пусть $L_{m}-$ д. о. в $L_{2}\left(\Omega_{n}\right)$ и



\[

L_{1} u=\sum_{|\alpha| \leqslant m-1} d_{\alpha}(x) D^{\alpha} u,

\]



тогда система собственных и присоединенных функций для д. о. $L_{m}+L_{1}$ полна в $L_{2}\left(\Omega_{n}\right)$. Однако разложение функции в ряд по этой системе, вообще говоря, расходится и суммируется со скобками по обобщенному методу Абеля.



Если область $\Omega_{n}$ неограничена, то для выполнения условий теоремы Кіелдыша надо налагать дополнительные услов́ия на рост коэффициентных функций д. $о$.



Мало изучены несамосопряженные д. о. с непрерывной частью спектра. Это связано с тем, что для них нет аналога теоремы о спектральном разложении. Исключение составляет д. о., порожденный выражением $-\frac{d^{2} x}{d x^{2}}+q(x)$ при $x \in[0, \infty)$ или $x \in(-\infty, \infty)$ с комплекснозначной функцией $q(x)$. Пусть $\varphi(x, k)$ является решением уравнения $-u^{(4)}+q(x) u=k^{2} u$ при $0 \leqslant x<\infty$ и удовлетворяет начальным условиям $\varphi(0, k)=1$,





\end{document}